Definition-complete means a valid raw real; equivalence-complete means a
formal chain of \code{RealRaw.Equiv} theorems.  The table is a compact
release-style coverage suite, not a numerator for a calculus-completeness
percentage.  Its five rows are the deliberately distinct bridges formalized
by the \code{PiProofs.PiCoverageBridge} type and its \code{equivalent}
theorem.

The suite asks whether a construction useful elsewhere in calculus agrees
with an independently constructed value of the same constant.  Thus a row is
evidence for a capability gate, rather than a claim that five successful pi
proofs make five equal-sized contributions to scientific or engineering
calculus.  In particular, a continuous Peano--Baker theorem and a proof that
the selected exponential satisfies \(e^x{}'=e^x\) are major advances even
though neither creates a new pi row.

\begin{center}
\begin{tabular}{p{0.21\textwidth}p{0.29\textwidth}p{0.27\textwidth}cc}
\textbf{Capability exercised} & \textbf{Witnessed equivalence} &
\textbf{What it regression-tests} & \textbf{Def.} & \textbf{Equiv.}\\
\hline
Finite Archimedean geometry &
cross-fan circumference \(\sim\) area loop &
exact rational circle fans, square-root enclosures, and nested interval bounds &
\(\checkmark\) & \(\checkmark\)\\
Arctangent and power series &
geometric arctangent \(\sim\) Leibniz series &
circular inverse geometry against the alternating arctangent series &
\(\checkmark\) & \(\checkmark\)\\
Finite definite integration &
rectangle arctangent integral \(\sim\) Leibniz series &
certified rectangle integration and its endpoint bridge &
\(\checkmark\) & \(\checkmark\)\\
Compactified improper integration &
Cauchy full-line integral \(\sim\) area loop &
reciprocal-tail transport from a full line to a finite interval &
\(\checkmark\) & \(\checkmark\)\\
Algebraic-kernel integration &
reciprocal-quartic integral \(\sim\) Cauchy integral &
a nontrivial rational kernel, finite quadrature, and its integral identity &
\(\checkmark\) & \(\checkmark\)\\
\end{tabular}
\end{center}

The full \code{PiProofs.PiPresentation} registry retains supplementary
checked evaluators as regression implementations: polygon area, the two
perimeter normalizations, Nilakantha, and the single Machin power-series
formula.  They remain useful executable comparisons, but are not additional
readiness cells because they reuse one of the bridges above.

The direct \code{piCircumference} evaluator remains a geometric diagnostic:
its unresolved lower-endpoint refinement is a local square-root-enclosure
problem, not a missing calculus gate.  Lean now reduces it to
\code{AdjacentChordCurvatureMarginCoversFineWidths}: an explicit rational
margin comparing the two fine curvature certificates, their bisection widths,
and the coarse squared chord.  The reduction from that margin to the original
endpoint refinement is checked; the universal margin inequality remains open.
The next candidate for a new coverage cell is integration by parts:
\[
  \pi=4\int_0^1\arctan(x)\,\dd x+2\log 2.
\]
Its finite rational rectangle identity is now checked, but the cell remains
unmarked until the product/derivative/FTC bridge, common increasing/decreasing
piece refinement, and logarithm endpoint identity are all constructive.  The
arcsine/Newton and Gaussian rows remain future probes for inverse/square-root
integration and the exponential/full-line layers; advanced routes are
intentionally omitted from this progress table.
